% !TeX program = pdflatex
% Samodzielny podglad rysunku uzytego w Uwadze 26.1.5.
% Kompilowac z katalogu glownego projektu (geometria i topologia rozniczkowa).
\documentclass[11pt,a4paper]{article}
\usepackage[T1]{fontenc}
\usepackage[utf8]{inputenc}
\usepackage[polish]{babel}
\usepackage{lmodern}
\usepackage[margin=2.5cm]{geometry}
\usepackage{amsmath,amssymb}
\usepackage{tikz}
\usetikzlibrary{arrows.meta}
\pagestyle{empty}
\begin{document}
\begin{center}
  {\Large\bfseries Rogata sfera Alexandera: stadia skończone}\par
  \medskip
  \input{tex/rysunki/rogata-sfera-alexandera}
\end{center}

\noindent
Każdy narysowany etap jest gładką sferą po wygładzeniu nasad rogów.
Rogi są wypukłościami kuli, a nie dodatkowymi $1$-uchwytami.
Przecięcia na rzucie oznaczają przejścia nad i pod; powierzchnia
w przestrzeni nie przecina samej siebie. Właściwa rogata sfera
powstaje dopiero przez nieskończone powtarzanie tego samego
sprzężenia w coraz mniejszej skali.
\end{document}
